\section*{Introducción}

Las termocuplas son sensores de temperatura ampliamente utilizados en la industria y la investigación debido a su robustez, amplio rango de medición y capacidad para operar en entornos hostiles. Su funcionamiento se basa en el efecto Seebeck, un fenómeno físico por el cual se genera una diferencia de potencial eléctrico (tensión) en un circuito compuesto por dos metales o aleaciones diferentes cuyas uniones se mantienen a temperaturas distintas. Esta tensión, minúscula pero medible, es aproximadamente proporcional a la diferencia de temperatura entre la unión de medida (o junta caliente) y la unión de referencia (junta fría), lo que permite determinar la temperatura en puntos críticos de procesos industriales como el control de hornos, la monitorización de turbinas, la automatización de procesos químicos o en aplicaciones aeroespaciales. Su simplicidad constructiva, que carece de electrónica compleja en el punto de medición, las hace particularmente valiosas para medir temperaturas extremadamente altas o en lugares de difícil acceso.

